\documentclass[10pt,a4paper]{article}
\usepackage[utf8]{inputenc}
\usepackage{amsmath}
\usepackage{amsfonts}
\usepackage{amssymb}
\usepackage{enumerate}
\title{GSF-3100 \\
02-Mathématiques financières des obligations \\
Série d'exercices 1}

\begin{document}
\maketitle

\begin{itemize}
\item Dans ce document vous avez des questions à choix multiples allant de a) à d).  
\item Il y a seulement un choix possible par question.
\item Ce questionnaire contient des exercices en lien avec le principe de valeur temporelle de l'argent permettant de déterminer le prix d'une obligation
\item Contient 11 questions. 
\end{itemize}
 

\newpage

\section*{Exercice 1}
Trouver le taux effectif semestriel équivalent à un taux nominal annuel de 9 \% capitalisé en continu.

\begin{enumerate}[a)]
\item 4,65\%
\item 4,55\%
\item 4,60\%
\item 4,70\%
\end{enumerate}


\section*{Exercice 2}
Trouver le taux nominal quotidien capitalisé annuellement équivalent à un taux nominal bisannuel de 20 \% capitalisé trimestriellement.

\begin{enumerate}[a)]
\item 0,045\%
\item 0,015\%
\item 0,022\%
\item 0,028\%
\end{enumerate}


\section*{Exercice 3}
Trouver la valeur actuelle d’une annuité de 35 paiements hebdomadaires de 3 \$ si le taux nominal annuel capitalisé hebdomadairement est de 12 \%.

\begin{enumerate}[a)]
\item 100.35\$
\item 100.79\$
\item 101.40\$
\item 100.75\$
\end{enumerate}


\section*{Exercice 4}
Trouver la valeur future d’une annuité de 2 paiements semestriels de 5 \$ en supposant que le taux nominal semestriel capitalisé trimestriellement est de 7 \% et que le premier paiement aura lieu aujourd'hui (début de période).

\begin{enumerate}[a)]
\item 11.09\$
\item 10.99\$
\item 12.09\$
\item 11.40\$
\end{enumerate}


\section*{Exercice 5}
Trouver le prix d’une obligation d’une valeur nominale de 1000 \$ venant à échéance dans 21,5 ans et ayant un taux de coupon (TC) de 24 \% annuel payable quatre fois par année, soit à chaque trimestre, en supposant un taux nominal annuel capitalisé mensuellement de 12 \%.

\begin{enumerate}[a)]
\item 1000\$
\item 1495.35\$
\item 1910.59\$
\item 1904.9\$
\end{enumerate}

\section*{Exercice 6}
Trouver le prix d’une obligation perpétuelle (à vie) ayant un taux de coupon de 16 \% annuel payable deux fois par année et une valeur nominale de 1000 \$ en supposant un taux eectif semestriel de 6 \%.

\begin{enumerate}[a)]
\item 1433.3\$
\item 11333.3\$
\item 1336.3\$
\item 1400.3\$
\end{enumerate}


\section*{Exercice 7}
Trouver le prix $P_k$ d’une obligation d’une valeur nominale de 800 \$ venant à échéance dans 5,5 ans et ayant un taux de coupon de 10 \% annuel payable deux fois par année en supposant que le taux nominal annuel capitalisé semestriellement est de 12 \%.  Cette obligation se transigeait initialement à $P_0$ = 736,9050 \$ et le dernier coupon a été versé le 31 décembre dernier il y a exactement 126 jours. (Utiliser la méthode exact/exact.)

\begin{enumerate}[a)]
\item 787.1\$
\item 763.1\$
\item 767.1\$
\item 796.9\$
\end{enumerate}

\newpage
\section*{Exercice 8}
Trouver la cote d'une obligation ayant un prix $P_k$ = 85 \$ entre deux dates de coupon et une fraction de période $k = 0,87$ écoulée depuis le versement du dernier coupon d'une valeur $C = 4$ \$.

\begin{enumerate}[a)]
\item 81.00\$
\item 81.13\$
\item 84.00\$
\item 81.52\$
\end{enumerate}

\section*{Exercice 9}
À l'aide d'une calculatrice financière, trouver le taux de rendement à l’échéance eectif semestriel d’une obligation d’une valeur nominale de 10 000 \$ venant à échéance dans 5 ans et ayant un taux de coupon semestriel de 2,5 \% payable deux fois par année et un prix de 8980 \$.

\begin{enumerate}[a)]
\item 3.4\%
\item 3.5\%
\item 3.6\%
\item 3.7\%
\end{enumerate}



\section*{Exercice 10}
Trouver le taux de rendement courant d'une obligation ayant un taux de coupon annuel de 10 \% payable deux fois par année, une valeur nominale de 1000 \$ ainsi qu'un prix au marché de $P_0$  = 1100 \$.

\begin{enumerate}[a)]
\item 9,01\%
\item 9,99\%
\item 9,09\%
\item 9,49\%
\end{enumerate}

\newpage

\section*{Exercice 11}
Un investisseur possède une obligation payant, en deux versements, un coupon annuel de 12 \%.  L’obligation vient à échéance dans 10 ans et promet un taux de rendement annuel nominal de 10 \% (tous les taux nominaux sont capitalisés semestriellement).  L’investisseur a un horizon de 2 ans et prévoit réinvestir ses coupons à un taux annuel nominal de 8 \% pendant cette période.  Quel rendement effectif annuel l’investisseur prévoit-il réaliser s’il croit être en mesure de vendre son obligation à la fin de son horizon à un taux de rendement annuel nominal de 12 \%?

\begin{enumerate}[a)]
\item 0,66\%
\item 0,56\%
\item 0,60\%
\item 0,50\%
\end{enumerate}
\end{document}